\section{Exact Positive Support Under Gradient Clipping}
\label{sec:appendix-support}

This appendix states the one-batch reward-influence results summarized in
Section~\ref{sec:reward-geometry}.  Theorem~\ref{thm:reward-halfspace-support}
is proved here; Proposition~\ref{prop:normalized-reward-influence} is proved
with the other deferred results in
Appendix~\ref{sec:appendix-proofs}.

\subsection{Normalized One-Batch Proposal Set}

This subsection connects the reward-record attacker to the release halfspaces
of Proposition~\ref{prop:parameter-gate}.  For a
fixed on-policy batch of length $T$, let $J\in\mathbb R^{T\times d}$ contain
the policy-score row vectors, let $r\in\mathbb R^T$ be the logged rewards, let
$H_{ij}=\mathbf 1[j\ge i]$ be the reward-to-go operator, and let
$C=I-T^{-1}\mathbf 1\mathbf 1^\top$ be the centering matrix.  An attacker
choosing $\delta\in[-\epsilon,\epsilon]^T$ produces the centered return
vector
\begin{equation}
y(\delta)=CH(r+\delta).
\label{eq:centered-return}
\end{equation}
The fixed-batch assumption is appropriate for a reward-record-only attacker:
changing a reward after collection does not change the states, actions, or
score rows already in the batch.

\begin{proposition}[Standardized reward-influence set]
\label{prop:normalized-reward-influence}
Suppose the learner standardizes the centered reward-to-go advantages with
population standard deviation whenever $\|y(\delta)\|_2>0$.  Before optimizer
clipping, a gradient-ascent proposal of size $\alpha$ is
\begin{equation}
\widetilde\theta^+(\delta)=\theta+
\frac{\alpha}{\sqrt T}J^\top
\frac{y(\delta)}{\|y(\delta)\|_2}.
\label{eq:normalized-update}
\end{equation}
Consequently, the raw one-batch proposal set is the image of the centered
reward-to-go zonotope under normalization and the policy-score map; in general,
it is not a parameter-space zonotope.
\end{proposition}

Normalization is why a bounded reward box does not induce a bounded
parameter-space zonotope, and it is also what makes the influence question
tractable.  The support characterization below applies to one fixed batch;
later batches change $J$, $r$, and the visited states.  The implementation
skips standardization below a numerical tolerance, so each reported cone
witness must remain above that threshold.

\subsection{Exact Positive Support}

\begin{theorem}[Positive support under global gradient clipping]
\label{thm:reward-halfspace-support}
Let $B=J^\top/\sqrt T$, let the learner cap the global gradient norm at
$G>0$, and assume coordinatewise parameter clipping is inactive over the
reward box.  The post-cap update before parameter clipping is
\begin{equation}
 \theta_G^+(y)=\theta+
 \frac{\alpha By}{\max\{\|y\|_2,\|By\|_2/G\}}.
\end{equation}
For a release halfspace $a^\top P\theta\le b$, define
$q=\alpha B^\top P^\top a$, $y_0=CHr$, and $M=CH$.  For any fixed $t>0$, a
reward perturbation whose gate-coordinate increment is at least $t$ exists if
and only if the following second-order-cone problem is feasible:
\begin{equation}
\begin{split}
 \lambda&\ge0,\qquad
 -\epsilon\lambda\mathbf1\le v\le\epsilon\lambda\mathbf1,\\
 \bar y&=\lambda y_0+Mv,\qquad q^\top\bar y=1,\\
 t\|\bar y\|_2&\le1,\qquad
 (t/G)\|B\bar y\|_2\le1.
\end{split}
\label{eq:soc-support-feasibility}
\end{equation}
Feasibility is monotone in $t$, so bisection on
$[0,\min\{\|q\|_2,\alpha G\|P^\top a\|_2\}]$ computes the exact positive
support
\begin{equation}
\max\left\{0,\sup_{\|\delta\|_\infty\le\epsilon}
\frac{q^\top y(\delta)}
{\max\{\|y(\delta)\|_2,\|By(\delta)\|_2/G\}}\right\}.
\end{equation}
The construction remains valid when the centered-return zonotope contains the
origin; if $G=\infty$, the second norm constraint is omitted.  If this support
does not exceed the current halfspace margin, no admissible reward edit can
cross that halfspace in the batch.
\end{theorem}

\begin{proof}
Global norm clipping gives the stated maximum in the denominator.  The ratio
is positively homogeneous, so optimization over the zonotope
$Z=\{y_0+M\delta:\|\delta\|_\infty\le\epsilon\}$ is equivalent to optimization
over its conic hull.  The first line of~\eqref{eq:soc-support-feasibility} is
the perspective representation of that hull: $v=\lambda\delta$.  This is the
standard homogenization used in fractional and conic
optimization~\cite{schaible_fractional_duality_1976,bauschke_homogenization_2023}.

If a positive ratio reaches $t$, rescale its conic representative until
$q^\top\bar y=1$; its two denominator branches are then exactly the last line
of~\eqref{eq:soc-support-feasibility}.  Conversely, any feasible $\bar y$ is
nonzero because of the unit-numerator equality and recovers a reward-box ray
whose ratio reaches $t$.  Thus the origin cannot create the degenerate
feasibility present in a direct zonotope formulation.  Each fixed-$t$ problem
is a second-order feasibility-cone problem~\cite{belloni_soc_feasibility_2008};
monotonicity permits quasiconvex bisection~\cite{agrawal_dqcp_2020}.
Cauchy--Schwarz supplies both stated upper bounds.  The margin result follows
by adding the maximum increment to the current gate coordinate.
\end{proof}
